\documentclass[10.5pt,a4paper]{article}
\usepackage[T1]{fontenc}
\usepackage[utf8]{inputenc}
\usepackage[english,italian]{babel}
\usepackage{lmodern,microtype}
\usepackage[a4paper,margin=1.85cm,headheight=15pt]{geometry}
\usepackage{enumitem,tabularx,array,longtable,booktabs}
\usepackage{xcolor,hyperref,xurl,fancyhdr,titlesec,framed}
\definecolor{ddtadark}{RGB}{28,38,52}
\definecolor{ddtablue}{RGB}{42,88,135}
\definecolor{ddtagray}{RGB}{244,246,248}
\definecolor{ddtagreen}{RGB}{48,111,78}
\hypersetup{colorlinks=true,linkcolor=ddtablue,urlcolor=ddtablue,
  pdftitle={DDTA modello scheda di lettura v2 - LaTeX companion},
  pdfauthor={Documentation-Driven Threat Analysis research workspace}}
\pagestyle{fancy}\fancyhf{}\lhead{DDTA - scheda di lettura}\rhead{v2 LaTeX companion}\cfoot{\thepage}
\titleformat{\section}{\large\bfseries\color{ddtadark}}{}{0pt}{}
\setlist[itemize]{leftmargin=1.4em,itemsep=0.15em,topsep=0.2em}
\setlength{\parindent}{0pt}\setlength{\parskip}{0.28em}\setlength{\emergencystretch}{2em}
\newcommand{\blankline}{\par\noindent\rule{\linewidth}{0.25pt}\vspace{0.42em}}
\newcommand{\blanklines}[1]{\par\begingroup\count255=0\loop\ifnum\count255<#1\blankline\advance\count255 by 1\repeat\endgroup}
\newenvironment{promptbox}[1]{%
  \def\FrameCommand{\fcolorbox{ddtablue}{ddtagray}}%
  \MakeFramed{\advance\hsize-2\width\FrameRestore}\noindent\textbf{#1}\par\smallskip}{\endMakeFramed}
\newenvironment{notearea}[1]{\textbf{#1}\par\smallskip}{\par\medskip}
\begin{document}

\begin{center}
{\LARGE\bfseries DDTA - Modello scheda di lettura}\par
\vspace{0.2em}
{\large Versione 2 - LaTeX companion}\par
\vspace{0.35em}
{\small Ausilio di lettura. Il record citabile resta \texttt{notes/<SRC>.md} + \texttt{excerpts/<SRC>.excerpts.yml}.}
\end{center}
\vspace{0.4em}

\begin{promptbox}{Identita bibliografica e accesso}
\begin{tabularx}{\linewidth}{>{\bfseries}p{3.3cm}X}
Source ID & \\
Citation key & \\
Titolo & \\
Autori / ente & \\
Anno / versione & \\
Tipo di fonte & \\
Identificatore stabile & \\
Landing URL & \\
PDF / full-text legale & \\
File locale & \\
Data retrieval & \\
SHA-256 locale & \\
Verification state & \\
\end{tabularx}
\end{promptbox}

\section{1. Problema di ricerca e contributo}
\textit{Quale problema affronta la fonte? Quale contributo dichiara? Qual e il livello di astrazione?}
\blanklines{7}

\section{2. Artefatti iniziali e rappresentazione del sistema}
\textit{Da quali artefatti parte? Quale rappresentazione assume, crea, estrae o deriva? Distinguere fatti sorgente, modelli intermedi, viste e output.}
\blanklines{9}

\section{3. Traceability, change, automation e ruolo umano}
\textit{Come collega fonti, requirements, decisioni/evidenze e cambiamento? Cosa e automatizzato e cosa richiede review/decisione umana?}
\blanklines{9}

\section{4. Findings, limiti e relazione con DDTA / overlay}
\textit{Che cosa dimostra davvero? Quali limiti dichiara? Quali limiti emergono per DDTA? Quale research gap rimane aperto?}
\blanklines{10}

\section{Metodo di ricerca e corpus}
\blanklines{6}

\section{Contributo e findings principali}
\blanklines{8}

\section{Limitazioni dichiarate dagli autori}
\blanklines{6}

\section{Validita temporale}
\begin{tabularx}{\linewidth}{>{\bfseries}p{4.1cm}X}
Contribution type & \\
Time sensitivity & \\
Temporal status & \\
Evidence window & \\
Elementi ancora applicabili & \\
Elementi potenzialmente obsoleti & \\
Corroborazione contemporanea & \\
Uso consentito in tesi & \\
\end{tabularx}
\vspace{0.4em}

\section{Relazione con Documentation-Driven Threat Analysis}
\begin{notearea}{Proposizioni direttamente supportate}\blanklines{5}\end{notearea}
\begin{notearea}{Tensioni o contraddizioni}\blanklines{4}\end{notearea}
\begin{notearea}{Research gap lasciato aperto}\blanklines{4}\end{notearea}
\begin{notearea}{Implicazioni per Base Analysis / overlay}\blanklines{5}\end{notearea}

\section{Candidate thesis claims}
\blanklines{6}

\section{Citation-ready locations}
\textit{Annotare pagina PDF, pagina stampata, sezione e tipo di uso: C = quotation verificata; P = paraphrase fedele; I = interpretazione del ricercatore.}
\begin{longtable}{p{1.1cm}p{1.3cm}p{1.8cm}p{3.1cm}p{7.2cm}}
\toprule
ID & Tipo & Pagina & Sezione & Proposizione / nota\\
\midrule
\endhead
 & & & & \\
\addlinespace[1.1em]
 & & & & \\
\addlinespace[1.1em]
 & & & & \\
\addlinespace[1.1em]
 & & & & \\
\addlinespace[1.1em]
 & & & & \\
\bottomrule
\end{longtable}

\section{Follow-up sources}
\textit{Separare riferimenti foundational, competing/alternative ed empirical/evaluation. Una citazione bibliografica non diventa automaticamente una fonte DDTA registrata.}
\blanklines{7}

\vfill
{\footnotesize\color{ddtadark}\textbf{Governance note.} Questa scheda e un worksheet. Le affermazioni usate nella tesi devono risolversi a una fonte registrata e a una posizione esatta nell'excerpt ledger. Quotazione, parafrasi e interpretazione devono restare separate.}
\end{document}
